\section{Conclusion}\label{sec:conclusion}

We have traced a deliberate path from the simplest computational unit—an LLM as a text-in, text-out function—to increasingly sophisticated multi-agent systems. This journey did not require changing the LLM itself. Instead, we progressively wrapped the core model in layers of control: first with agents that observe and act, then with execution models that orchestrated parallel and streaming invocations, next with sub-agent hierarchies that decomposed complex problems, and finally with orchestration patterns that coordinated multiple independent systems. At each step, the model remained constant while the architecture around it grew in complexity and capability.

This insight is fundamental: the LLM is an uncommitted tool. Its power comes not from its individual inference but from how we compose multiple inferences within a structured process. The four canonical patterns we examined—pipeline, fan-out/fan-in, master-slave, and feedback loop—are the vocabulary of this composition. These patterns are not LLM-specific inventions; they echo classical concurrent and distributed computing. Their importance here is that they are small enough to understand completely, composable enough to combine into larger systems, and general enough to apply across an enormous range of agent tasks. A system that orchestrates document analysis, code review, or scientific hypothesis testing may use these same patterns, differing only in the details of task decomposition and resource constraints.

Selecting the right architecture is not primarily a question of model capability. Given a task, the architecture emerges from the task's structure: its inherent parallelism, its feedback requirements, its hierarchical depth, and the latency constraints we impose. A well-designed agent system for real-time reasoning may use tight feedback loops and shallow hierarchies. A batch system for asynchronous analysis may embrace deep pipelines and wide fan-outs. Neither choice is superior; each is the natural fit to its problem. This decoupling between model selection and architectural design frees practitioners to reason about systems clearly.

As LLM-based applications move from research prototypes to production systems at scale, these patterns will become foundational. We can already see them emerging in production orchestration frameworks and multi-agent platforms. Understanding them deeply—not merely as abstract diagrams but as concrete execution models with resource, latency, and consistency implications—is essential for anyone building intelligent systems. The progression we have outlined, from simple agents to complex hierarchies, provides both a map and a toolkit for that future. The LLM remains at the heart of the system, but the architecture around it is where the real engineering begins.
